% Exemplo: texto em fórmulas, casos, operadores, conjuntos numéricos, teoremas e unidades
% Compilar: latexmk -pdf recursos.tex              Testado com: TeX Live 2024 ou mais recente
\documentclass[a4paper]{article}
\usepackage[T1]{fontenc}
\usepackage{lmodern}
\usepackage[brazilian]{babel}
%%BEGIN:recursospreambulo
\usepackage{amsmath, amssymb}   % amssymb: \mathbb e outros símbolos
\usepackage{amsthm}             % ambientes de teorema
\usepackage{siunitx}            % números e unidades
\sisetup{locale = DE, output-decimal-marker = {,}}
\newtheorem{teorema}{Teorema}
\theoremstyle{definition}
\newtheorem{definicao}{Definição}
\DeclareMathOperator{\sen}{sen}  % seno em português
%%END:recursospreambulo

\begin{document}

%%BEGIN:recursos
\begin{definicao}
O valor absoluto de $x \in \mathbb{R}$ é
\[
  \lvert x \rvert =
  \begin{cases}
    x,  & \text{se } x \geq 0, \\
    -x, & \text{caso contrário.}
  \end{cases}
\]
\end{definicao}

\begin{teorema}[Pitágoras]
Em um triângulo retângulo de catetos $a$, $b$ e hipotenusa $c$, vale $a^2 + b^2 = c^2$.
\end{teorema}

\begin{proof}
Fica como exercício. Note que $\sen^2\theta + \cos^2\theta = 1$.
\end{proof}

A velocidade medida foi de \qty{12,5}{\metre\per\second}, com erro de \qty{0,3}{\percent}.
%%END:recursos

\end{document}
